% ══════════════════════════════════════════════════════════════════════════════
% CHAPTER 7 — DISCUSSION  [v31 — revised for defence]
% ══════════════════════════════════════════════════════════════════════════════

This chapter interprets the empirical results from Chapter~\ref{ch:results}.
The purpose is not to repeat the tables, but to explain what the estimates imply
about central bank communication as a market signal, why the results differ
across methods and institutions, and where the evidence should be interpreted
cautiously. The discussion is organised around the two research questions. RQ1
asks whether hawkish--dovish tone explains announcement-window market reactions
after controlling for the contemporaneous rate surprise. RQ2 asks whether
central bank language predicts near-term policy outcomes and two-year yield
changes in an expanding-window setting.

\section{Summary and interpretation of key findings}
\label{sec:disc_summary}

The central conclusion is that central bank language contains measurable
information, but the information is not uniform across assets, institutions, or
measurement methods. Text is most informative when the dependent variable is
close to the monetary-policy object being communicated: short-maturity yields,
rate-decision classifications, and next-meeting policy moves. Text is weaker
when the outcome is more exposed to other forces, such as equity returns,
foreign exchange, long yields, and future two-year yield changes.

For RQ1, the evidence is concentrated rather than universal. In FOMC statements,
the clearest market-reaction result is the two-year Treasury yield. Both the
Random Forest and BERT scores are positively associated with UST 2Y
reactions in the controlled regressions, and the full horse race raises the
common-sample $R^2$ to 0.262 (adjusted $R^2 = 0.249$). This is economically
coherent. A two-year yield is closely linked to the expected near-term path of
policy, so it should be the asset most likely to respond to language about
future tightening or easing. The result is robust to Newey-West standard errors
and survives fixed-effect controls. By contrast, the FOMC equity and
exchange-rate regressions show little evidence that text adds meaningful
explanatory power beyond the rate surprise.

For ECB statements, the strongest RQ1 result is the Euro Stoxx 50. The Random
Forest score is positively associated with SX5E reactions ($p=0.028$), and the
coefficient remains significant in the three-way horse race ($p=0.045$). The
positive sign is consistent with an information-effect interpretation: language
associated with tighter policy is also language associated with stronger
macroeconomic conditions, and equity markets respond positively to that
information. However, this result becomes marginal ($p=0.102$) when chair and
era fixed effects are added, suggesting that part of the association may
reflect systematic differences between ECB presidents (whose tenures coincide
with different rate regimes) rather than within-regime communication effects.
Two other ECB outcomes---OIS 10Y and EUR/USD---have negative adjusted $R^2$ in
all specifications, meaning the text scores add no explanatory power beyond
the rate-surprise control after adjusting for degrees of freedom.

The press-conference evidence reinforces the interpretation that communication
is information-rich, but the mapping from tone to market reaction is less
stable. In the FOMC press-conference sample ($N=91$), the Random Forest score
is significantly associated with EUR/USD and Treasury yields. However, the
yield signs are negative---the opposite of the standard prediction that
hawkish language raises yields. This result is not driven by outliers
(removing the two most influential meetings strengthens the coefficient) and
is robust to Newey-West standard errors. It suggests that post-2011 press
conferences combine discussion of tightening conditions with risk-management
reassurance, so the RF score captures a mixed signal rather than a pure
hawkish yield shock. For ECB press conferences, BERT is the clearest
equity signal, while the Random Forest contributes to EUR/USD in the horse
race. The broader conclusion is that press conferences are information-rich,
but the type of information is more heterogeneous than in statements.

The dual-mandate decomposition adds a second layer to the interpretation. For
FOMC statements, decomposing the dictionary score into policy, inflation,
employment, and residual-risk components improves $R^2$ modestly. For ECB
statements, the gains are larger, especially for OIS\_MP2 (7.0 percentage
points) and OIS 10Y (8.1 percentage points). For OIS 10Y, a partial $F$-test
confirms that the additional sub-category regressors are jointly significant
($p=0.027$), and the adjusted $R^2$ rises from $-0.002$ under the aggregate
score to $0.055$ under the full decomposition. This is not evidence that the
ECB has a Federal Reserve-style dual mandate. Rather, it shows that ECB market
reactions are component-specific: policy-action language, inflation language,
growth language, and residual-risk language do not have the same
financial-market implications.

For RQ2, the evidence is clearer. Both Random Forest and BERT forecasts
produce positive direct expanding-window $R^2$ values for future rate-decision
targets in all eight corpus--method pairs. BERT performs especially
strongly for rate decisions, reaching direct out-of-sample $R^2$ values of
0.457 for FOMC statements and 0.281 for ECB press conferences at the
next-meeting horizon. By contrast, the same methods do not forecast future
two-year yield changes: all sixteen yield forecasts (two methods $\times$ four
corpora $\times$ two horizons) deliver negative direct out-of-sample $R^2$.
The conclusion is therefore unambiguous: central bank language helps predict
subsequent policy actions but does not predict subsequent yield changes.

This distinction is not contradictory. A document can contain information about
what the central bank is likely to do at the next meeting while still failing to
forecast the yield change realised between meetings. Markets may price much of
the expected policy path immediately at the announcement, leaving little
predictable movement for the subsequent yield window. Future two-year yields also
respond to incoming inflation releases, labour-market data, risk premia, and
global financial shocks that occur after the communication event. The RQ2 result
therefore separates policy predictability from market predictability: language
helps forecast the central bank's next action, but not the full path of market
repricing that occurs after the meeting.

These findings imply a layered view of communication. Dictionary scores provide
transparent measures of explicit stance. Random Forest scores capture language
that is predictive of realised rate-decision signs. BERT scores capture
sentence-level semantic tone and are particularly useful in forecasting future
rate decisions. The three methods are complementary, but they should not be
interpreted as measuring exactly the same economic object.

\section{Economic interpretation: policy news and information news}
\label{sec:disc_policy_information}

The results are easiest to interpret through the distinction between policy news
and information news. Policy news refers to communication that changes market
expectations about the future policy path. Information news refers to
communication that changes market beliefs about the state of the economy. The
same hawkish sentence can contain both. For example, language about persistent
inflation may signal future tightening, which is negative for equities through
higher discount rates. But it may also reveal that demand remains strong, which
can be positive for expected cash flows. The observed asset-price response
depends on which channel dominates.

This ambiguity explains why the equity results are not uniformly negative. For
FOMC statements, the text coefficients for equities are weak. This is consistent
with offsetting channels: hawkish language may raise expected rates while also
signalling stronger economic conditions. For ECB statements, the positive SX5E
response to the Random Forest score is consistent with the information channel
being stronger in that sample. This is plausible because ECB communications
often combine policy-action language with an assessment of euro-area conditions.
If markets interpret tighter-policy language as evidence that the Governing
Council sees the economy as resilient enough to absorb tightening, equities can
rise rather than fall. However, this interpretation should be treated with
caution. In the present framework, the two channels cannot be separated
structurally---a negative equity coefficient could equally be attributed to the
discount-rate channel, and a positive one to the information channel, making the
interpretation unfalsifiable without additional identifying restrictions such as
the sign-restriction approach of \citet{jarocinski2020}.

The fixed-effect and forward/current results provide partial evidence for the
dual-channel view. For ECB statements, forward-looking tone is negatively
associated with SX5E returns, while current-condition tone is positively
associated with SX5E returns. The signs are economically intuitive. Hawkish
forward-looking language points to future tightening and is negative for
equities. Hawkish current assessment points to stronger current economic
conditions and is positive for equities. The aggregate score partly masks this
difference because it pools the two channels. While this pattern is harder to
explain with pure post-hoc rationalisation than the single-coefficient results,
it still falls short of a formal structural identification.

For yields, the expected sign under both the standard monetary-policy channel
and the information effect is positive: hawkish language should raise yields by
shifting the expected policy path upward. This prediction holds cleanly in FOMC
statements, where the two-year Treasury response is the strongest yield result
and the ten-year response is weaker. However, the prediction fails in FOMC press
conferences, where the RF score has a significantly negative yield coefficient.
This sign reversal means that the yield channel is not as ``unambiguous'' as the
methodology section's theoretical framework suggests. The negative sign is
consistent with post-2011 press conferences mixing tightening rationale with
risk-management reassurance, but it could also reflect a period-specific pattern
related to the zero lower bound, QE tapering, and rapid tightening eras. Long
yields include term premia, long-run inflation expectations, and global risk
factors; a simple hawkish--dovish score should therefore not be expected to
explain them as cleanly.

\section{RQ2 and the distinction between policy predictability and market predictability}
\label{sec:disc_rq2_interpretation}

The RQ2 results draw an important distinction between predicting policy decisions
and predicting market prices. Future rate decisions are partly under the control
of the central bank and are shaped by the same reaction function that appears in
its communication. If the central bank repeatedly describes inflation as too
high, policy as insufficiently restrictive, or labour markets as tight, that
language can predict subsequent rate hikes because both the language and the
future decision reflect the same internal assessment.

One interpretive caveat applies to the strength of this conclusion. All three
text scores are highly autocorrelated: the first-order autocorrelation of
\netscore{} is 0.87 in FOMC statements, and the RF and BERT scores range from
0.55 to 0.78 across corpora. Part of the rate-decision forecast power may
therefore reflect the text score acting as a slow-moving proxy for the policy
stance rather than capturing meeting-specific forward-guidance signals. The
incremental $R^2$ analysis partially addresses this concern by controlling for
the current rate decision, and the RF text forecast remains significant beyond
inertia ($p=0.005$ for FOMC statements). But the high autocorrelation means
the text score carries more information about the general policy trajectory
than about the specific incremental signal at any given meeting.

Future two-year yield changes are different. They are market outcomes between
meetings. They depend on new macroeconomic data, risk sentiment, term premia,
liquidity conditions, geopolitical shocks, and revisions in the expected policy
path that may occur after the original communication. Even if a statement helps
forecast the next policy decision, it may not forecast the subsequent yield
change because the market may already price the decision at the announcement, or
because new information arrives before the next meeting.

This explains the central RQ2 pattern. Direct out-of-sample $R^2$ values are
positive for future rate decisions in all eight corpus--method pairs but negative
for all sixteen yield forecasts. Some regression-based incremental $R^2$ values
for yield targets are positive, especially for BERT in ECB
specifications, but those gains should be interpreted cautiously because the
direct forecast benchmark fails in every case. The conclusion is that central
bank language predicts the central bank's subsequent policy choices but not the
market's future repricing.

The contrast between Random Forest and BERT also matters. The Random
Forest is trained on document-level bigram features and captures phrase-level
patterns associated with policy actions. BERT classifies sentences and
therefore captures more local semantic context. The fact that BERT
performs better in many rate-decision forecasting specifications suggests that
sentence-level tone is useful for extracting forward-looking policy information.
This is particularly visible in press conferences, where central bankers explain
the conditional path of policy more explicitly than in short statements.

\section{FOMC versus ECB}
\label{sec:disc_fomc_ecb}

The results differ across the FOMC and the ECB in ways that are consistent with
their institutional structures. FOMC statements are committee-drafted documents
and are closely watched for changes in wording. This makes them suitable for
short-maturity yield reactions: small changes in policy-path language can affect
two-year Treasury yields. The FOMC statement results therefore fit the idea that
textual tone matters most for assets closely tied to the expected funds-rate
path.

ECB statements and press conferences have a different structure. The ECB press
conference has existed for the entire sample, while FOMC press conferences begin
only in 2011. ECB communication therefore provides a longer press-conference
history across more regimes. ECB policy communication also contains distinctive
terms around price stability, financial fragmentation, transmission, negative
rates, and asset purchases. These institutional features make ECB communication
particularly sensitive to the information channel. The positive SX5E response
to ECB statement RF tone is consistent with markets reading some hawkish ECB
language as macroeconomic news rather than only as monetary tightening, though
the loss of significance under chair fixed effects ($p=0.102$) means that this
interpretation rests partly on cross-president variation.

The ECB rate-decision benchmark also differs from the FOMC benchmark. The FOMC
rate decision is relatively straightforward because changes are usually recorded
in 25-basis-point increments. The ECB has multiple policy rates, and during the
negative-rate period the deposit facility rate became the binding marginal rate.
The thesis uses a consistent benchmark for classification and forecasting, but
some ECB meetings include non-standard movements or changes in the corridor
structure. This may weaken the relationship between text labels and the
rate-decision sign for the ECB, especially in periods when the policy stance is
transmitted through instruments other than the main refinancing rate.

\section{Dictionary, Random Forest, and BERT: trade-offs}
\label{sec:disc_method_tradeoffs}

The method comparison is not a ranking from worst to best. Each method measures
a different object and carries a different identification risk. The dictionary is
most transparent, the Random Forest is most directly linked to the realised
rate-decision sign, and BERT is most context-sensitive at the sentence
level. The empirical question is therefore not which score is universally best,
but which score is most informative for a given outcome and how much additional
complexity is justified.

The three scoring methods differ in transparency, flexibility, and context
sensitivity. The dictionary is the most transparent. Its categories are known in
advance, every score can be traced to individual term matches, and the
sub-scores are easy to interpret. This makes the dictionary useful for economic
explanation and for the dual-mandate decomposition. Its weakness is rigidity.
A fixed term list cannot adapt automatically to new communication styles,
changing policy instruments, or subtle contextual shifts.

The Random Forest is more flexible. It learns corpus-specific bigrams that
predict rate-decision signs, and it can identify institutional phrasings that a
hand-built dictionary misses. However, this advantage comes with two specific
interpretive costs. First, because the Random Forest is trained on realised rate
decisions, its high RF--rate agreement (84--95\% across corpora) is a training
convergence diagnostic, not an independent validation of economic stance.
Second, feature inspection reveals that approximately 18\% of total feature
importance in FOMC statements and 20\% in ECB statements is attributable to
leakage-adjacent terms---speaker names, voting-roll phrases, and institutional
boilerplate that correlate with specific eras rather than with economic stance.
The remaining $\sim$80\% of importance is carried by economically interpretable
terms, but the leakage share means the \rfscore{} is partly influenced by
era-specific vocabulary. The RF score is therefore best interpreted as a
decision-predictive communication measure, not as a pure structural measure of
hawkishness.

The horse-race regressions that include all three scores simultaneously face a
collinearity issue. The pairwise correlation between the dictionary and
BERT scores is 0.81 in FOMC statements, producing variance inflation
factors around 3.0 for these two variables. While below the conventional concern
threshold of 5--10, this collinearity means that individual coefficient
significance in the horse race is unreliable. When a score loses significance
in the joint model, it may reflect absorbed information rather than genuine
redundancy. The $R^2$ comparison across single-score models is therefore more
informative than individual horse-race coefficients.

BERT offers the strongest context sensitivity. It processes sentences
rather than isolated n-grams and can in principle distinguish between phrases
that a dictionary treats similarly. The results show that BERT is
especially useful in RQ2 rate-decision forecasting. At the same time, its
discrete classification is heavily skewed: 81 of 91 scored FOMC press
conferences are labelled dovish. While the continuous \bertscore{} retains
variation within the dovish range (the regression results are driven by this
continuous variation, not the discrete labels), the systematic dovish tilt
reflects domain transfer from the FOMC-trained sentence classifier to broader
or ECB-specific language. Transformer scores should therefore be treated as a
useful benchmark, not as an unquestioned ground truth. The sentence-level
classifier is also applied once to the full corpus without retraining in the
RQ2 expanding-window design; only the regression layer mapping \bertscore{} to
forecast targets is fit in an expanding window. This means the sentence
predictions themselves are not strictly out-of-sample, although the economic
mapping is.

The thesis therefore does not identify a single dominant method. Instead, the
methods answer different questions. The dictionary is best for transparent
measurement and decomposition. The Random Forest is best for identifying
language associated with rate decisions but carries a leakage risk that future
work should address. BERT is strongest for sentence-level semantic
information and near-term rate-decision forecasting, but its dovish shift and
domain-transfer limitations should be kept in mind. The three-way comparison is
valuable because it reveals when these objects coincide and when they diverge.

\section{Comparison with prior literature}
\label{sec:disc_literature}

The results are consistent with the broader literature showing that central bank
communication contains information beyond the current rate decision, but they
also qualify that conclusion. The event-study literature argues that monetary
policy announcements include both a target component and a path component. The
empirical design of this thesis follows that logic by controlling for the
rate-surprise component and testing whether text explains residual market
movements. The strongest evidence appears at the short end of the yield curve
for the FOMC and in ECB equities, rather than uniformly across all asset classes.
This suggests that text-based tone captures only part of the path and
information content embedded in announcements.

The results also connect to dictionary-based studies of monetary-policy tone.
Like earlier dictionary work, the thesis finds that explicit hawkish and dovish
language tracks monetary-policy cycles and can be related to market outcomes.
However, the evidence here is more conditional than a simple narrative would
suggest. A single aggregate dictionary score is not uniformly significant across
markets. The decomposition results show why: different components of tone can
have opposite signs. This supports the view that monetary-policy text should not
always be collapsed into one index without checking whether policy-action,
inflation, growth, and risk language carry different signals.

The machine-learning results align with work showing that supervised models can
extract meaningful structure from central bank language. The Random Forest's
high OOB and balanced accuracy confirms that bigram features contain strong
information about the contemporaneous decision class. But the holdout results
(which drop by up to 23 percentage points relative to OOB) and the
feature-inspection evidence (18--20\% leakage-adjacent importance) also show
the limits of this approach. Central bank language changes across regimes, and
some predictive features are institutional rather than purely economic. The
blocklist reduces this problem but cannot remove it completely.

The BERT results are consistent with the idea that contextual language
models add value when the relevant signal is sentence-level and semantic rather
than phrase-level and fixed. This is most visible in RQ2 forecasting. Still, the
model's systematic dovish tilt in some corpora and the cross-institutional
transfer issue mean that transformer scores should be treated as a benchmark,
not as an unquestioned ground truth.

The null and weak results are also informative. The absence of robust effects in
many FX, equity, long-yield, breakeven, and sovereign-spread specifications---and
the negative adjusted $R^2$ for ECB OIS 10Y and EUR/USD---suggests that
hawkish--dovish tone is not a sufficient statistic for all market news released
around central bank meetings. This is an important boundary condition for the
communication literature. It indicates that text measures are most useful when
the outcome is close to the policy path, and less reliable when the asset price
embeds several competing risk premia or macro-financial channels.

\section{Limitations}
\label{sec:disc_limitations}

Several limitations should be emphasised.

\emph{Textual stance versus market surprise.}
A central limitation is that textual hawkishness and market hawkishness
are not identical. The dictionary, Random Forest, and BERT measures
capture the absolute restrictiveness of communication. Financial markets,
however, react to surprises relative to prior expectations. This creates
a mechanical problem for text-based measurement: meetings with rate hikes
naturally contain more restrictive language, even when the hike was fully
expected. Conversely, a realised hike can be interpreted as dovish if
markets had expected a larger increase or stronger forward guidance.

The thesis addresses this issue by controlling for high-frequency
rate-surprise measures in the event-study regressions and by comparing
text labels to both realised rate-decision labels and market-surprise
labels. Nevertheless, the text scores should be interpreted as measures
of communication stance, not as fully expectation-adjusted structural
monetary-policy shocks.

\textbf{identification issue 2 paragraphs above added KMC 
}
\emph{Constructed measures.}
The text scores are constructed measures, not observed economic variables. The
dictionary depends on term selection and threshold choices; the Random Forest
depends on the training target, the feature-selection procedure, and the
blocklist; BERT depends on sentence-level model predictions and the
domain of the training data. The use of three methods reduces dependence on any
single measurement approach, but it does not eliminate measurement error.

\emph{Descriptive Random Forest and residual leakage.}
The Random Forest classification exercise is descriptive. Because the RF is
trained to classify the contemporaneous rate-decision sign, high RF--rate
agreement should not be interpreted as an independent discovery that language
contains policy information. Moreover, approximately 18\% of total feature
importance in FOMC statements and 20\% in ECB statements is attributable to
leakage-adjacent terms that correlate with specific speakers or eras rather than
with economic stance. While the remaining $\sim$80\% is economically
interpretable, the leakage share means that the \rfscore{} used in RQ1
regressions is partly influenced by era-specific vocabulary. A rerun with the
extended blocklist would be needed to confirm that the market-reaction results
are robust to this contamination.

\emph{Endogeneity of the RF score in RQ1.}
The RF score is trained on the contemporaneous rate decision. Controlling for
MP1 removes the surprise component, but the expected policy action is embedded
in both the RF score and the market reaction. The RF coefficient in the
market-reaction regressions may therefore partly reflect the mechanical
relationship between expected policy and asset prices, rather than a pure
communication effect. The dictionary and BERT scores, which are not
trained on the rate decision, provide a cleaner test of the communication
channel. The fact that BERT shows similar significance for FOMC UST 2Y
($p=0.020$) supports the interpretation that the result is not purely a
mechanical RF artefact.

\emph{ECB policy-rate benchmark.}
ECB policy has been conducted through several instruments, especially during the
negative-rate and asset-purchase periods. A single rate-decision sign cannot
fully capture changes in the deposit facility rate, the main refinancing rate,
asset-purchase programmes, forward guidance, or transmission-protection tools.
This is likely most relevant for ECB observations after 2014.

\emph{Uneven sample sizes.}
FOMC statements cover a long period, but FOMC press conferences begin only in
2011 ($N=91$). The FOMC press-conference forecast evaluation window is
particularly short (as few as 53 evaluated forecasts for RF rate-decision
targets). ECB press conferences cover a longer period, but some long-maturity
OIS variables are available only for much smaller samples ($N=113$ for OIS 10Y).
These sample constraints limit the statistical power and generalisability of the
press-conference and long-maturity results.

\emph{Multiple testing.}
The analysis reports results across four corpora, multiple outcomes, three
scoring methods, and several regression variants (horse race, fixed-effect
ladder, forward/current decomposition, dual-mandate decomposition), generating
in excess of 200 individual specifications. No formal multiple-testing
correction has been applied. While the patterns that are consistent across
methods---such as the positive rate-decision forecasting result, which holds
for all eight corpus--method pairs---are less likely to reflect spurious
significance, individually starred coefficients should be interpreted with
appropriate caution.

\emph{Text score persistence.}
All three text scores are highly autocorrelated ($\rho_1 \approx 0.55$--$0.87$
across corpora). Part of the RQ2 forecasting power may therefore reflect the
text score acting as a slow-moving proxy for the overall policy stance rather
than capturing meeting-specific forward-guidance information. The incremental
$R^2$ analysis controls for the current rate decision, which partially addresses
this concern, but a single-lag inertia control may not fully absorb the
slow-moving stance information embedded in the text scores.

\emph{BERT look-ahead and domain transfer.}
The sentence-level BERT classifier is applied once to the full corpus and is
not retrained in the expanding-window design. Only the regression layer mapping
\bertscore{} to forecast targets is fit on past data. If the classifier was
fine-tuned on an external labelled dataset, this is analogous to applying a
fixed dictionary---the predictions do not learn from the thesis corpus. However,
the distinction should be stated explicitly. The classifier also exhibits a
systematic dovish shift, particularly in press conferences (81 of 91 FOMC press
conferences labelled dovish), likely reflecting domain transfer from
FOMC-trained models.

\emph{Event-study associations, not causal effects.}
The event-study design identifies announcement-window associations, not fully
structural causal effects. Narrow windows reduce the probability of confounding
news, but they cannot prove that every price movement within the window is
caused only by the text score. Controlling for the rate surprise helps isolate
residual communication content, but the coefficient should still be read as an
association between measured tone and market reaction.

\emph{Forecast evaluation.}
The expanding-window design prevents look-ahead bias, yet it also reduces the
amount of training data available early in the sample. The direct out-of-sample
$R^2$ values are reported as point estimates without confidence intervals; a
Clark-West or bootstrap test would provide a formal significance assessment.
Negative direct out-of-sample $R^2$ values for yield targets should not be read
as evidence that communication is irrelevant for yields at the announcement;
rather, they show that communication is not sufficient to forecast subsequent
yield changes over the next one or two meetings.

\section{Implications}
\label{sec:disc_implications}

The first implication is methodological. Researchers should be cautious about
using a single text score as if it measured one stable object. Central bank
communication contains policy intentions, economic assessments, risk framing,
and institutional language. These components can have different and sometimes
opposite market effects. Decomposition and method comparison are therefore not
optional robustness checks; they are central to interpretation.

The second implication concerns central bank communication itself. Markets do
not only react to whether language is hawkish or dovish. They react to why it is
hawkish or dovish. Hawkishness based on current economic strength can raise
equities, while hawkishness based on future tightening can depress them. This
means that central banks cannot assume that tighter language will always be
interpreted as contractionary news. The market response depends on whether the
communication is read as policy news or information news.

The third implication concerns forecasting. Central bank language can help
predict future policy decisions, but it does not predict future yield movements.
This distinction matters for practitioners. A text model may be useful for
anticipating policy actions without being profitable or accurate as a
stand-alone yield-forecasting tool.

The fourth implication is a measurement caution. Supervised text classifiers
trained on policy outcomes can achieve high in-sample agreement with the
training target, but feature leakage (institutional boilerplate, speaker names,
era-specific phrases) can inflate apparent performance and contaminate
downstream regression estimates. Feature-importance diagnostics and blocklists
are therefore necessary complements to accuracy metrics, not optional
appendix material.

\section{Directions for future research}
\label{sec:disc_future_research}

Future research could extend this thesis in several directions. First, ECB
policy measurement could be improved by constructing a multi-instrument
decision benchmark that combines the MRO, deposit facility rate, asset-purchase
announcements, and forward-guidance changes. This would better reflect the ECB's
post-2014 operating framework.

Second, the dictionary could be made dynamic. Instead of applying one fixed
vocabulary across all regimes, future work could estimate era-specific or
institution-specific dictionaries while preserving the transparent category
structure. This would test whether the meaning of stance-bearing language
changes across conventional, zero-lower-bound, negative-rate, and post-COVID
periods.

Third, the BERT component could be retrained or fine-tuned on ECB
sentences. The current design is useful as a benchmark, but an ECB-specific
sentence classifier would provide a cleaner test of whether transformer-based
tone improves on the dictionary and Random Forest for euro-area communication.
Retraining the classifier within the expanding-window design would also
eliminate the look-ahead concern noted in Section~\ref{sec:disc_limitations}.

Fourth, the Random Forest could be rerun with the extended blocklist identified
in Appendix~\ref{app:a3_rf_features}. This would quantify how much the 18--20\%
leakage-adjacent feature importance affects the RQ1 regression coefficients and
the RQ2 expanding-window forecasts. A before-and-after comparison would test
whether the market-reaction results are robust to removing era-specific
vocabulary.

Fifth, future work could separate press-conference prepared remarks from Q\&A
answers. The current cleaning removes journalist questions and keeps central
bank answers, but the economic content of prepared remarks and spontaneous
answers may differ. A separate decomposition could test whether markets react
more strongly to scripted or unscripted communication.

Sixth, the policy-news and information-news channels could be separated
structurally. The current framework identifies a reduced-form association
between text and asset prices but cannot distinguish the two channels.
Combining text scores with the sign-restriction or heteroskedasticity-based
identification strategies used in the structural monetary-policy literature
would provide a more rigorous test of whether hawkish language operates through
the discount-rate channel, the information channel, or both.

Seventh, the out-of-sample $R^2$ values could be tested formally using
Clark-West or Diebold-Mariano tests. The current evaluation reports point
estimates without confidence intervals. Bootstrapping the expanding-window
procedure and reporting confidence bands would strengthen the RQ2 conclusions.

Eighth, future research could combine text with voice, facial expression, or
speaker-level uncertainty measures in press conferences. The present thesis
focuses on written language. Press conferences, however, also communicate tone
through delivery. Combining textual and non-verbal channels could provide a more
complete measure of central bank communication.